% =====================================================================
%  CARNET D'ENTRAÎNEMENT — CHIMIE — FICHE 3 : Noyau atomique, atome
%  THÈME 1 : Particules subatomiques (charge, masse, énergie cinétique)
%  FICHE TUTO — à lire avant les exercices
% =====================================================================
\documentclass[11pt,a4paper]{article}

\usepackage[utf8]{inputenc}
\usepackage[T1]{fontenc}
\usepackage{lmodern}
\usepackage[french]{babel}

\usepackage{amsmath,amssymb,mathtools}
\usepackage{xcolor}
\usepackage{colortbl}
\usepackage{tikz}
\usetikzlibrary{arrows.meta,positioning,calc,shadings}
\usepackage[most]{tcolorbox}
\usepackage{enumitem}
\usepackage{booktabs}
\usepackage{array}
\usepackage[a4paper,margin=2.1cm,headheight=26pt,headsep=9pt]{geometry}
\usepackage{fancyhdr}
\usepackage{titlesec}
\titlespacing*{\section}{0pt}{6pt}{3pt}
\titleformat{\section}{\large\bfseries\color{primarydark}}{}{0pt}{}
\usepackage[hidelinks]{hyperref}

% ---------- Palette ----------
\definecolor{primary}{RGB}{124,78,140}
\definecolor{primarydark}{RGB}{74,44,92}
\definecolor{defcol}{RGB}{0,114,178}
\definecolor{thmcol}{RGB}{0,140,120}
\definecolor{formulacol}{RGB}{198,93,9}
\definecolor{attncol}{RGB}{200,40,55}
\definecolor{protoncol}{RGB}{205,50,50}
\definecolor{neutroncol}{RGB}{105,108,120}
\definecolor{electroncol}{RGB}{25,95,205}

% ---------- Notation nucléaire ----------
\newcommand{\nuc}[3]{\prescript{#1}{#2}{\mathrm{#3}}}
\newcommand{\pp}{p^{+}}
\newcommand{\ee}{e^{-}}
\newcommand{\nz}{n^{0}}

% ---------- En-tête ----------
\pagestyle{fancy}
\fancyhf{}
\fancyhead[L]{\small\textcolor{primarydark}{\textbf{Fiche 3 — Thème 1}} : Particules subatomiques}
\fancyhead[R]{\small\textcolor{primarydark}{\textsc{Tuto}}}
\fancyfoot[C]{\small\thepage}
\renewcommand{\headrulewidth}{0.8pt}
\renewcommand{\headrule}{\hbox to\headwidth{\color{primary}\leaders\hrule height \headrulewidth\hfill}}

% ---------- Boîtes ----------
\tcbset{
  boxbase/.style={enhanced,breakable,boxrule=0.9pt,arc=2.5pt,
     left=3mm,right=3mm,top=2.2mm,bottom=2.2mm,
     fonttitle=\bfseries,coltitle=white,toptitle=0.5mm,bottomtitle=0.5mm},
}
\newtcolorbox{cle}[1][]{boxbase,colback=thmcol!6,colframe=thmcol,
  title={\textsc{Point clé}\ifx\relax#1\relax\relax\else\textmd{ : \itshape #1}\fi}}
\newtcolorbox{formule}[1][]{boxbase,colback=formulacol!7,colframe=formulacol,
  title={\textsc{Formule}\ifx\relax#1\relax\relax\else\textmd{ : \itshape #1}\fi}}
\newtcolorbox{piege}[1][]{boxbase,colback=attncol!7,colframe=attncol,
  title={\textsc{Piège à éviter}\ifx\relax#1\relax\relax\else\textmd{ : \itshape #1}\fi}}
\newtcolorbox{ex}[1][]{boxbase,colback=defcol!5,colframe=defcol,
  title={\textsc{Exemple}\ifx\relax#1\relax\relax\else\textmd{ : \itshape #1}\fi}}

\setlist{leftmargin=1.5em,topsep=2pt,itemsep=1.5pt}

\begin{document}

% ---------- Bandeau titre ----------
\begin{tikzpicture}
\node[rectangle,rounded corners=6pt,fill=primary,text=white,
      minimum width=\textwidth,minimum height=1.5cm,align=center,inner sep=6pt]
  {\Large\bfseries Thème 1 — Les particules subatomiques\\[2pt]
   \normalsize\mdseries charge, masse et énergie cinétique};
\end{tikzpicture}

\vspace{3pt}
{\small\itshape À lire avant les exercices. Objectif : reconnaître les trois particules par leur
\textbf{charge} et leur \textbf{masse}, calculer la \textbf{charge d'un ensemble} de particules, et
relier une énergie cinétique d'électron à sa vitesse.}

\section*{1. Trois particules, deux grandeurs}

Deux particules vivent dans le \textbf{noyau} (proton, neutron), la troisième peuple le \textbf{nuage
électronique} (électron). On les caractérise par leur \textbf{charge} et leur \textbf{masse}.

\begin{center}
\renewcommand{\arraystretch}{1.3}
\begin{tabular}{l c c c}
\toprule
\rowcolor{primary!14}
\textbf{Particule} & \textbf{Où ?} & \textbf{Charge} & \textbf{Masse (approx.)} \\
\midrule
Proton $\ \pp$    & noyau              & $+e = +1{,}6\times10^{-19}$ C & $\approx 1{,}7\times10^{-27}$ kg \\
Neutron $\ \nz$   & noyau              & $0$ (neutre)                  & $\approx 1{,}7\times10^{-27}$ kg \\
Électron $\ \ee$  & nuage électronique & $-e = -1{,}6\times10^{-19}$ C & $\approx 9{,}1\times10^{-31}$ kg \\
\bottomrule
\end{tabular}
\end{center}

\noindent Deux faits à graver dans le marbre :
\begin{itemize}
  \item le \textbf{proton} et le \textbf{neutron} ont \emph{quasiment la même masse} : ce sont eux qui
        font presque toute la masse de l'atome ;
  \item l'\textbf{électron} est \emph{environ $2000$ fois plus léger} que le proton — sa masse n'est
        \emph{pas} nulle, mais elle est négligeable devant celle d'un nucléon.
\end{itemize}

\begin{formule}[charge élémentaire et rapport des masses]
\[ e = 1{,}6\times10^{-19}\ \text{C}
   \qquad\qquad
   \frac{m(\ee)}{m(\pp)} = \frac{9{,}1\times10^{-31}}{1{,}7\times10^{-27}} \approx 0{,}0005 \approx \frac{1}{2000}. \]
La charge élémentaire $e$ est la plus petite charge « à l'état libre ». Le proton porte $+e$,
l'électron $-e$ (charges \emph{opposées}, même valeur absolue), le neutron est neutre.
\end{formule}

\begin{cle}[l'atome est surtout fait de vide]
Le noyau mesure $\sim 10^{-15}$ m ; l'atome entier $\sim 10^{-10}$ m, soit $\sim 10^{5}$ fois plus grand.
Comme le volume croît comme le \emph{cube} de la taille, l'espace « habité » par les électrons est
$\sim(10^{5})^{3}=10^{15}$ fois plus grand que le noyau. Le constituant principal, \textbf{en volume},
d'un atome est donc le \textbf{vide}.
\end{cle}

\section*{2. Charge d'un ensemble de particules}

La charge d'un paquet de particules identiques est simplement le \textbf{nombre de particules} multiplié
par la charge de \emph{chacune}. C'est une proportionnalité.

\begin{formule}[charge d'un ensemble]
\[ \boxed{\,Q = N \times q_{\text{unité}}\,}
   \qquad\text{par ex. }\ Q(\text{$N$ protons}) = N \times (+1{,}6\times10^{-19}\ \text{C}). \]
Pour des électrons, on prend $q_{\text{unité}} = -1{,}6\times10^{-19}$ C (charge négative).
\end{formule}

\begin{ex}[charge de $10^{12}$ protons]
Chaque proton porte $+1{,}6\times10^{-19}$ C. Pour $N=10^{12}$ protons :
\[ Q = 10^{12}\times(+1{,}6\times10^{-19}) = +1{,}6\times10^{\,12-19} = +1{,}6\times10^{-7}\ \text{C}. \]
\textbf{Astuce des exposants :} on multiplie les mantisses ($10^{12}\times1{,}6$) et on \emph{additionne}
les exposants ($12+(-19)=-7$).
\end{ex}

\section*{3. Énergie cinétique d'un électron (pour aller plus loin)}

À l'échelle atomique on exprime l'énergie en \textbf{électron-volts} (eV) : il faut la convertir en
\textbf{joules} avant tout calcul avec une masse en kg.

\begin{formule}[énergie cinétique et vitesse]
\[ \boxed{\,E_c = \tfrac{1}{2}\,m\,v^{2}\,}
   \qquad\Longleftrightarrow\qquad
   \boxed{\,v = \sqrt{\dfrac{2E_c}{m}}\,} \]
$E_c$ en joules (J), $m$ en kg, $v$ en m/s. Conversion utilisée dans le livre :
$1\ \text{eV} = 1{,}66\times10^{-19}$ J (la valeur physique exacte est $1{,}602\times10^{-19}$ J).
\end{formule}

\begin{ex}[méthode en 3 temps]
Un électron ($m=9{,}11\times10^{-31}$ kg) a $E_c = 18{,}22$ eV. Sa vitesse ?
\begin{enumerate}
  \item \textbf{Convertir} : $E_c = 18{,}22\times1{,}66\times10^{-19}$ J.
  \item \textbf{Isoler} : $v=\sqrt{2E_c/m}=\sqrt{\dfrac{2\times18{,}22\times1{,}66\times10^{-19}}{9{,}11\times10^{-31}}}$.
  \item \textbf{Simplifier} : $2\times18{,}22=36{,}44$ et $36{,}44/9{,}11=4$ exactement ;
        $10^{-19}/10^{-31}=10^{12}$. Donc
        $v=\sqrt{4\times1{,}66\times10^{12}}=2\sqrt{1{,}66}\times10^{6}\approx 2{,}58\times10^{6}$ m/s.
\end{enumerate}
\end{ex}

\begin{piege}
\textbf{Ne dites jamais que l'électron « n'a pas de masse ».} Il est très léger ($\sim2000\times$ moins
que le proton), mais sa masse n'est pas nulle. \; \textbf{N'oubliez jamais la conversion eV~$\to$~J} :
un $E_c$ laissé en eV avec une masse en kg donne un résultat faux de plusieurs ordres de grandeur.
\end{piege}

\end{document}
